% s23_hyp_optional: the hypotheses of the optional extension (W2), (R_W), (FA^K); input from sec/s23.tex
% (moved here from the inputs section; labels unchanged)

\begin{hypothesis}[\Wtwo: the Wick sector at arbitrary amplitudes]\label{hyp:s03:W2}
Let $a$ be an admissible scheme. (a) For every finite family of real $2$-forms and real functions, every $J\ge J_0(f,h)$ and \emph{every} real coefficient vector $(z,w)$, the limit $\lim_kZ^a_{k,J}(z,w)$ exists and lies in $(0,\infty)$. (b) For every $q\in C^\infty_c(\R^3)$ with $q\ge0$,
\[
\sup_{J\ge J_q}\ \lim_{k\to\infty}Z^{a}_{k,J}(0,q)<\infty ,
\]
where $J_q$ is the least $J$ with $\supp q$ at $\ell^\infty$-distance greater than $1$ from $\R^3\setminus\CJ$.
\end{hypothesis}

\begin{remark}[on \Wtwo]\label{rem:s03:W2}
(i) Clause (a) is \cref{hyp:s03:W1} with the threshold on the positive scalar part removed. On the cone of amplitudes $\sup q\le s_1<\sigma_1(M_1)$, clause (b) is proved in this paper in $L^1$ form (\cref{cor:s11:W2prime}, on the comparison convention; on $\NDK$ in $\limsup$ form by \cref{rem:s04:F}). Above the threshold both clauses are open; the exact obstruction is the free-energy bound $\Bminus{u}$ at amplitude $u$ above the threshold (\cref{hyp:s23:Bminus} below): the $\eps$-uniform stability of King's torus model with its bare coefficient shifted by $-2u$.
(ii) No clause of \cref{thm:s13:main} uses \Wtwo. The Wick result of the main theorem is Form~(B) on the tube $\Dtube$, where only the proved cone instance enters, through \cref{thm:s11:W}; \Wtwo\ is a hypothesis of Form~(A) (\cref{thm:s23:FormA}) and of \RW. In the rotation part, \Wtwo\ for a scheme $a$ is the $\mathrm W$-admissibility $\Wadm{a}$ of \cref{rot:hyp:adm}.
\end{remark}

\begin{hypothesis}[\RW: rotation invariance of the Wick sector at arbitrary amplitudes]\label{hyp:s03:RW}
For an admissible scheme $a$ satisfying \Wtwo, the full object $Z^a_\infty=\lim_J\lim_kZ^a_{k,J}$ of the comparison family (whose existence and entireness follow from \Wtwo\ as in \cref{lem:s08:W1}) is invariant under every Euclidean motion for all complex coefficients, $F$ transforming as a $2$-form and $W^a$ as a scalar.
\end{hypothesis}

\begin{remark}[on \RW]\label{rem:s03:RW}
\RW\ enters only in the Euclidean clause of Form~(A) (\cref{thm:s23:FormA}). Form~(B) has no Euclidean clause and does not assume it. The rotation part proves the curvature sector; its Wick-sector theorems (\cref{rot:thm:gen,rot:thm:closure}) give the $O(3)$ clause of \RW\ for a scheme $a$ (its translations are \cref{thm:s23:FormA}(ii)) under the torus comparison with scalar sources and the closing hypothesis for that scheme (\cref{rot:hyp:Ka,rot:hyp:TIa}).
\end{remark}

\begin{hypothesis}[\FAK: Form~(A) on King's convention]\label{hyp:s03:FAK}
On King's convention $\NDK$, under \Wtwo\ read for $\NDK$: the limit functional $Z^{a,K}_\infty$ is entire in all coefficients; the lattice symmetries $\Oh$ about every point hold for the joint field at all amplitudes; reflection positivity holds at arbitrary complex coefficients; the measure exists off the tube; and the Euclidean clause of the Wick sector holds for the King-family object.
\end{hypothesis}

\begin{remark}[on \FAK]\label{rem:s03:FAK}
\FAK\ is the King-convention counterpart of Form~(A) as the comparison-convention results give it (\cref{thm:s23:FormA}, last paragraph). The main theorem makes no use of it: its Wick leg on King's convention is Form~(B) on the tube (\cref{thm:s13:TWK}, \cref{prop:s13:WK}).
\end{remark}
